\documentclass[10pt,a4paper]{amsart}
\usepackage[margin=19mm]{geometry}
\usepackage{amsmath,amssymb,mathtools}
\usepackage[scr=rsfs,cal=euler]{mathalfa}
\usepackage{xfrac}
\usepackage[unicode,hidelinks,bookmarks=false]{hyperref}
\newcommand{\ie}{i.e.\ }
\newcommand{\cf}{cf.\ }
\newcommand{\resp}{respectively\ }
\newcommand{\Subsection}{\subsection}
\providecommand{\nobreakdash}{\nobreak\hbox{-}}
\setlength{\emergencystretch}{2em}
\multlinegap0pt
\usepackage{amssymb}
\usepackage[shortlabels]{enumitem}
\newtheorem{lemma}{Lemma}
\title{The Reduced Smith Group of a Kneser Graph and Its Application to Group-Valued Magic Maps}
\author{Ahmet Batal}
\date{Source: arXiv:2609.12546v1 (CC BY 4.0)}
\begin{document}
\maketitle

\noindent\emph{Source and licence.} The Reduced Smith Group of a Kneser Graph and Its Application to Group-Valued Magic Maps. Version arXiv:2609.12546v1 (CC BY 4.0), distributed by arXiv under the Creative Commons Attribution 4.0 International licence, http://creativecommons.org/licenses/by/4.0/, as recorded in the arXiv licence metadata for this exact version and shown on the abstract page https://arxiv.org/abs/2609.12546v1. Full licence notice: \texttt{licenses/arxiv-2609.12546-CC-BY-4.0.txt}.\par\smallskip

\noindent\textit{Editorial scope.} Counted unit: the article's lemma lem:bier-mobius (source0.tex lines 392--394). The article's complete original proof of it is delivered with it (source0.tex lines 396--451, one \texttt{proof} environment of 56 lines). No mathematics is written, repaired or re-derived: the statement and the proof are the article's own lines, in the article's own notation, and no retained line is substituted or re-flowed.  No further source-local context is needed: every cross-reference of every retained line resolves inside the two delivered spans.  Nothing was omitted from any retained span, so the package carries no omission record.  No retained line contains a citation, so the wrapper carries no bibliography and no bibliography entry.  No retained line contains a figure environment, an image-inclusion command or drawing code, so no image is delivered and the package declares no asset dependency.  Editorial formatting only, in addition: an AMS \texttt{amsart} preamble that restates the article's own declarations and macros used by the retained lines, namely the environments \texttt{lemma} on separate counters, together with the standard packages those lines need.  The retained environments print in this document as Lemma 1 in order of appearance, whereas the article prints the same items with the numbering of its own full document.  The complete native source of the article as distributed in the arXiv e-print for this exact version is bundled unchanged as \texttt{sources/210127/source0.tex}.
\begin{lemma}\label{lem:bier-mobius}
The matrix $Q$ is unimodular.
\end{lemma}

\begin{proof}
Define an endomorphism $H$ of $\mathbb Z^{\mathcal B}$ by
\[
  H_{\alpha,\beta}
  =(-1)^{|\alpha|}[\alpha\subseteq\beta],
  \qquad \alpha,\beta\in\mathcal B.
\]
We first show that $Q=PH$.
Fix $u\in V$ and $\beta\in\mathcal B$. Expanding the corresponding matrix
entry gives
\begin{equation}\label{eq:ph-entry}
  (PH)_{u,\beta}
  =
  \sum_{\alpha\in\mathcal B}
  [\alpha\subseteq u](-1)^{|\alpha|}[\alpha\subseteq\beta]
  =
  \sum_{\substack{\alpha\in\mathcal B\\
                   \alpha\subseteq u\cap\beta}}
  (-1)^{|\alpha|}
  =
  \sum_{\alpha\subseteq u\cap\beta}
  (-1)^{|\alpha|}.
\end{equation}
The last equality follows from the fact that $\mathcal B$ is closed under
taking subsets. Indeed, write $\beta=\{b_1<\cdots<b_j\}$. If
$\alpha=\{b_{i_1}<\cdots<b_{i_t}\}\subseteq\beta$, then
$b_{i_s}\ge2i_s\ge2s$ for every $1\le s\le t$. Thus $\alpha$ is
standard and hence belongs to $\mathcal B$. The last sum in
\eqref{eq:ph-entry} equals $1$ if $u\cap\beta=\emptyset$. If
$u\cap\beta\ne\emptyset$, then $u\cap\beta$ has equally many subsets of even
and odd cardinality. Indeed, toggling the membership of any fixed element of
$u\cap\beta$ gives a bijection between these two classes of subsets. Hence
the last term of \eqref{eq:ph-entry} is $
  [u\cap\beta=\emptyset]
  =Q_{u,\beta}.
$
Therefore $Q=PH$.

We next show that $H$ is unimodular. If
$H_{\alpha,\beta}\ne0$, then $\alpha\subseteq\beta$ and
$|\alpha|\le|\beta|$. Hence $H$ is block upper triangular with respect to the
level order on $\mathcal B$. If $\alpha$ and $\beta$ have the same
cardinality and $\alpha\subseteq\beta$, then $\alpha=\beta$. The diagonal
block indexed by $\mathcal B_j$ is therefore
\[
  (-1)^jI_{m_j},
\]
where $I_{m_j}$ denotes the $m_j\times m_j$ identity matrix.
It follows that
\[
  \det H=\prod_{j=0}^r(-1)^{j m_j}\in\{1,-1\},
\]
and hence $H$ is unimodular. The identity
$Q=PH$ now shows that $Q$ is unimodular
because both $P$ and $H$ are unimodular.
\end{proof}

\end{document}
